%
% vox-music-player-transport.tex
%
% Z Specification for the voxd Music Player's phase-3 transport controls.
%
% Phase 1/2 (docs/vox-music-player.tex) modelled a two-mode PlayerView --- idle
% and playing --- projecting voxd's single active source onto a lux scene, with
% three invariants:
%
%   I1  at most one album playing
%   I2  now-playing present iff playing
%   I3  a played album is catalogued
%
% Phase 3 adds a transport bar (prev / play-pause / next / stop) and a real
% pause/resume capability. The player becomes a THREE-mode machine --- idle,
% playing, paused --- and the phase-2 invariants generalise. This model pins the
% phase-3 PlayerView and the seven invariants T1--T7 of
% docs/vox-music-player-transport.md across the transport transitions.
%
% It supersedes the two-mode PlayerView of phase 2 by forward integration: the
% same object, one extra mode. It does not contradict phase 2 --- every phase-2
% behaviour survives with `paused' added. The brief once proposed a single
% departure (T4 restricting play to the idle mode, against the phase-2
% start-or-switch PlayerPlay), which this model surfaced as Fork A; the leader
% resolved it to SWITCH, keeping the phase-2 semantics with no regression, so
% PlayAlbum carries no mode guard. See Section~\ref{sec:forks}.
%

\documentclass[a4paper,10pt,fleqn]{article}
\usepackage[margin=1in]{geometry}
\usepackage{fuzz}
\usepackage[colorlinks=true,linkcolor=blue,citecolor=blue,urlcolor=blue]{hyperref}

\begin{document}

\title{The voxd Music Player, Phase 3: A Z Specification of the Transport}
\author{Formal model of the transport state machine}
\date{July 2026}
\hypersetup{
  pdftitle={The voxd Music Player, Phase 3: A Z Specification of the Transport},
  pdfauthor={Punt Labs},
  pdfsubject={Z Specification},
  pdfcreator={fuzz/probcli}
}
\maketitle

\section{Introduction}

The transport bar adds four buttons to the \texttt{vox.music} scene --- prev,
play/pause, next, stop --- and, with them, a genuine \emph{pause/resume}
capability the earlier player lacked. The phase-1/2 player
(\texttt{docs/vox-music-player.tex}) was a two-mode machine: \texttt{idle} or
\texttt{playing}. Phase~3 inserts a third mode, \texttt{paused}, between them: a
held source that is neither torn down (as \texttt{stop} would) nor progressing
(as \texttt{playing} does).

This specification pins the phase-3 \emph{PlayerView} --- the projection of
\texttt{voxd}'s single active source onto the scene --- and the seven invariants
\texttt{T1}--\texttt{T7} the transport must preserve across every transition. It
reuses the style, the naming, and the opaque \texttt{ALBUM} type of
\texttt{docs/vox-music-player.tex}, and it supersedes that document's two-mode
\texttt{PlayerView} by forward integration: the same object with one more mode,
not a parallel one. The transport carries no new playback machinery of its own;
\texttt{pause}/\texttt{resume}/\texttt{prev} name new \texttt{ProgramService}
capabilities whose mechanism (a \texttt{SIGSTOP}/\texttt{SIGCONT} of the playback
subprocess, or a playback-loop suspend) is the daemon specialist's decision. The
model fixes the behaviour and the invariants, not the mechanism.

The relationship to the phase-2 invariants is exact: \texttt{T1} is \texttt{I1}
widened to count \texttt{paused} as active; \texttt{T2} is \texttt{I2} widened so
the cursor is present in both active modes; \texttt{T7} is \texttt{I3},
unchanged. \texttt{T3}, \texttt{T4}, \texttt{T5}, and \texttt{T6} are new with the
transport.

\section{Basic Types, Modes, and Constants}
\label{sec:types}

An \texttt{ALBUM} is the opaque catalog identifier of
\texttt{docs/vox-music-player.tex} --- the id \texttt{vox music list} prints and
the scene's Play button carries. It is a catalog key, never a filesystem path;
the model treats it as opaque.

\begin{zed}
  [ALBUM]
\end{zed}

The player is now a three-mode machine: nothing playing, one album playing, or
one album held.

\begin{zed}
  PlayerMode ::= idle | playing | paused
\end{zed}

The play/pause button carries one of three glyphs. \texttt{pauseGlyph} is
\texttt{U+23F8} ($\Box\Box$, shown while \emph{playing} --- press it to pause);
\texttt{playGlyph} is \texttt{U+23F5} ($\triangleright$, shown while
\emph{paused} --- press it to resume); \texttt{inertGlyph} stands for the
disabled, greyed transport when \texttt{idle}. The glyph is a derived
observation of the mode, never an independent variable
(Section~\ref{sec:state}); \texttt{T6} pins the two live cases as
biconditionals.

\begin{zed}
  GLYPH ::= pauseGlyph | playGlyph | inertGlyph
\end{zed}

A \texttt{ZBOOL} is a Z boolean, named to avoid the ProB/B clash with
\texttt{BOOL}. It flags the derived \texttt{advancing} observation --- whether the
cursor is auto-progressing --- which is how \texttt{T3} enters the state
predicate as a fact rather than as prose.

\begin{zed}
  ZBOOL ::= ztrue | zfalse
\end{zed}

The part cursor and the album's part count are bounded so that ProB can explore
the reachable state space; \texttt{maxPart} is a model bound, not a product
limit. The phase-2 model wrote the cursor over \texttt{\nat}; here it is refined
to $0 \upto maxPart$ precisely so the exploration terminates. Three parts is
enough to exercise every boundary: \texttt{prev} at part~1, \texttt{next} at
part~$M$, and the interior.

\begin{axdef}
  maxPart : \nat_1
  \where
  maxPart = 3
\end{axdef}

\section{State}
\label{sec:state}

The \texttt{PlayerView} carries the catalog of saved albums; the mode; the
optional active album (a set of at most one, the album that is playing
\emph{or} paused); the now-playing cursor \texttt{(index, of)} --- the 1-based
part position and the album's part count, both zero when idle; and two derived
observations, \texttt{advancing} and \texttt{glyph}, that lift \texttt{T3} and
\texttt{T6} into the state predicate.

The predicate carries every invariant. Read the conjuncts in the order of the
brief's \texttt{T1}--\texttt{T7}.

\begin{schema}{PlayerView}
  catalogued : \power ALBUM \\
  mode : PlayerMode \\
  album : \power ALBUM \\
  index : 0 \upto maxPart \\
  of : 0 \upto maxPart \\
  advancing : ZBOOL \\
  glyph : GLYPH
  \where
  \# album \leq 1 \\
  (mode \in \{ playing, paused \} \iff \# album = 1) \\
  (mode = idle \iff album = \emptyset) \\
  (mode \in \{ playing, paused \} \implies (1 \leq index \land index \leq of \land of \geq 1)) \\
  (mode = idle \implies (index = 0 \land of = 0)) \\
  (advancing = ztrue \iff mode = playing) \\
  (glyph = pauseGlyph \iff mode = playing) \\
  (glyph = playGlyph \iff mode = paused) \\
  (glyph = inertGlyph \iff mode = idle) \\
  album \subseteq catalogued
\end{schema}

The conjuncts discharge the invariants as follows.

\begin{description}
  \item[\texttt{T1} --- single active source.] $\# album \leq 1$ with
    $(mode \in \{ playing, paused \} \iff \# album = 1)$: at most one album is
    active, and it is active in exactly the two non-idle modes. A \texttt{paused}
    album still occupies the one active slot, so a \texttt{play} cannot start a
    second alongside it.
  \item[\texttt{T2} --- now-playing iff active.] The cursor is live
    ($1 \leq index \leq of$, $of \geq 1$) exactly in the active modes, and is
    pinned to $(0,0)$ when \texttt{idle}. Widened from \texttt{I2} to admit
    \texttt{paused}: a paused album keeps its cursor, unlike a stopped one.
  \item[\texttt{T3} --- paused is suspended.] $advancing = ztrue \iff mode =
    playing$. Hence $mode = paused \implies advancing = zfalse$: the cursor
    auto-advances only while playing. \texttt{AutoAdvance}
    (Section~\ref{sec:autoadvance}) is guarded by $advancing = ztrue$, so it can
    never fire in the paused mode --- the state face of ``paused does not
    progress''.
  \item[\texttt{T4} --- transition guards.] Not a state invariant but a family
    of preconditions, one per operation; it is discharged in
    Sections~\ref{sec:play}--\ref{sec:autoadvance}, where each transition names
    the mode it is enabled from. \texttt{PlayAlbum} is the one transition
    enabled from \emph{every} mode --- it starts from \texttt{idle} and switches
    from \texttt{playing} or \texttt{paused} (Section~\ref{sec:forks},
    Fork~A) --- while \texttt{pause}, \texttt{resume}, \texttt{stop},
    \texttt{prev}, and \texttt{next} each carry the mode guard named in their
    schema. It cannot live in the state predicate, because it constrains
    \emph{steps}, not \emph{states}; recording it as prose here would be
    exactly the mistake the brief warns against, so it is recorded instead as
    the precondition line of every operation schema.
  \item[\texttt{T5} --- cursor bounds.] The active-mode conjunct fixes
    $1 \leq index \leq of$. The no-op boundary cases --- \texttt{prev} at part~1,
    \texttt{next} at part~$of$ --- are properties of those two operations
    (Sections~\ref{sec:prev}--\ref{sec:next}), stated there as total
    definitions.
  \item[\texttt{T6} --- glyph reflects state.] The three glyph biconditionals.
    The two the brief names --- $pauseGlyph \iff playing$, $playGlyph \iff
    paused$ --- plus the idle case that makes \texttt{glyph} total.
  \item[\texttt{T7} --- catalogued.] $album \subseteq catalogued$: the active
    album, playing or paused, is one the catalog holds. This is \texttt{I3}
    unchanged.
\end{description}

Because \texttt{PlayerMode} partitions into three disjoint inhabitants, the two
membership biconditionals are consistent: $mode = idle \iff album = \emptyset
\iff \# album = 0$, and $mode \neq idle \iff \# album = 1$. The derived
\texttt{advancing} and \texttt{glyph} are total functions of \texttt{mode}, so
every state fixes them uniquely; they add no freedom, only a name for what the
scene reads off the mode.

\section{Initialisation}

The daemon starts with some catalog and nothing playing: the idle shape, an
inert transport, a cursor at zero. The boot catalog is left unconstrained --- any
subset of \texttt{ALBUM} --- so that the animator explores every catalog shape,
from empty to full; an initialisation establishes a state, it does not consume an
input, which is why the phase-2 \texttt{cat0?} input is dropped here in favour of
a nondeterministic \texttt{catalogued'}.

\begin{schema}{PlayerInit}
  PlayerView'
  \where
  mode' = idle \\
  album' = \emptyset \\
  index' = 0 \\
  of' = 0 \\
  advancing' = zfalse \\
  glyph' = inertGlyph
\end{schema}

The \texttt{advancing'} and \texttt{glyph'} lines are forced by the state
invariant from $mode' = idle$; they are written out so the idle projection is
manifestly unique for the animator. Only \texttt{catalogued'} is free.

\section{Transport Transitions}

Each operation names, on its first predicate line, the mode it is enabled from
--- that line is \texttt{T4}. The framing lines then say what is held. The
derived \texttt{advancing'} and \texttt{glyph'} are left to the state invariant
of \texttt{PlayerView'}: each is a total function of \texttt{mode'}, so the
after-state fixes them, and writing them out in every schema would only repeat
the invariant.

\subsection{PlayAlbum --- start or switch to a saved album}
\label{sec:play}

A client plays the album \texttt{target?}, whose ready-part count is
\texttt{count?}. It fires from \emph{any} mode: from \texttt{idle} it
\emph{starts} a source; from \texttt{playing} or \texttt{paused} it
\emph{switches}, the newly chosen album displacing the active one and starting
at part~1. The only guard is that the album be catalogued and non-empty
($count? \geq 1$); there is no mode precondition, which is why the album-list
Play buttons stay enabled in every mode. The after-state is one album playing,
cursor at part~1, regardless of the before-mode.

This is the phase-2 \texttt{PlayerPlay} start-or-switch verb, carried into
phase~3 unchanged --- there is no mode guard, and no regression from phase~2.
The resolution of the fork this once raised is recorded in
Section~\ref{sec:forks}.

\begin{schema}{PlayAlbum}
  \Delta PlayerView \\
  target? : ALBUM \\
  count? : 1 \upto maxPart
  \where
  target? \in catalogued \\
  catalogued' = catalogued \\
  mode' = playing \\
  album' = \{ target? \} \\
  index' = 1 \\
  of' = count?
\end{schema}

Because $count? \geq 1$, the after-state satisfies the cursor-validity invariant
($1 \leq 1 \leq of'$, $of' \geq 1$); because $target? \in catalogued$,
\texttt{T7} holds. The after-mode is \texttt{playing}, so $advancing' = ztrue$
and $glyph' = pauseGlyph$ by the invariant.

\subsection{Pause --- suspend in place}
\label{sec:pause}

From \texttt{playing}, \texttt{pause} holds the source where it is: the mode
becomes \texttt{paused}, and the album and cursor are untouched. \texttt{T4}:
enabled only from \texttt{playing}.

\begin{schema}{Pause}
  \Delta PlayerView
  \where
  mode = playing \\
  catalogued' = catalogued \\
  mode' = paused \\
  album' = album \\
  index' = index \\
  of' = of
\end{schema}

The cursor is held exactly, so \texttt{T2} and \texttt{T5} carry over. The
after-mode is \texttt{paused}, so $advancing' = zfalse$ (\texttt{T3}: no further
auto-advance) and $glyph' = playGlyph$ (\texttt{T6}: the button now offers
resume).

\subsection{Resume --- continue a suspended album}
\label{sec:resume}

From \texttt{paused}, \texttt{resume} continues from where playback stopped: the
mode returns to \texttt{playing}, the cursor unchanged. \texttt{T4}: enabled only
from \texttt{paused}. \texttt{Resume} is the exact inverse of \texttt{Pause} on
the mode, the identity on everything else.

\begin{schema}{Resume}
  \Delta PlayerView
  \where
  mode = paused \\
  catalogued' = catalogued \\
  mode' = playing \\
  album' = album \\
  index' = index \\
  of' = of
\end{schema}

The after-mode is \texttt{playing}, so $advancing' = ztrue$ and $glyph' =
pauseGlyph$ again.

\subsection{Stop --- tear down, return to idle}
\label{sec:stop}

From either active mode --- \texttt{playing} or \texttt{paused} ---
\texttt{stop} tears the source down and returns the view to \texttt{idle}: no
album, no cursor, an inert transport. \texttt{T4}: enabled from
$\{ playing, paused \}$. The brief settles stop-while-paused explicitly: a
paused album, stopped, goes to \texttt{idle} exactly as a playing one does; there
is no ``paused-then-stopped'' residue.

\begin{schema}{Stop}
  \Delta PlayerView
  \where
  mode \in \{ playing, paused \} \\
  catalogued' = catalogued \\
  mode' = idle \\
  album' = \emptyset \\
  index' = 0 \\
  of' = 0
\end{schema}

Emptying the album empties the cursor, so \texttt{T2} is restored to its idle
shape; $advancing' = zfalse$ and $glyph' = inertGlyph$ follow from $mode' =
idle$.

\subsection{Prev --- previous part, floored at 1}
\label{sec:prev}

From either active mode, \texttt{prev} steps the cursor back one part. At part~1
it is a no-op --- the total definition the brief asks for: \texttt{prev} is
enabled in every active state, and its effect at the floor is the identity.
\texttt{T4}: enabled from $\{ playing, paused \}$. The mode is held: a paused
player stays paused, a playing player stays playing; \texttt{prev} moves the
cursor, it does not resume.

\begin{schema}{Prev}
  \Delta PlayerView
  \where
  mode \in \{ playing, paused \} \\
  catalogued' = catalogued \\
  mode' = mode \\
  album' = album \\
  of' = of \\
  (index > 1 \implies index' = index - 1) \\
  (index = 1 \implies index' = index)
\end{schema}

Since $index \geq 1$ in every active state, the two cases are exhaustive. When
$index > 1$, $index' = index - 1 \geq 1$; when $index = 1$, $index' = 1$; either
way $1 \leq index' \leq of$, so \texttt{T5} holds and the no-op at part~1 is
exactly the $index = 1$ case.

\subsection{Next --- next part, capped at M}
\label{sec:next}

From either active mode, \texttt{next} steps the cursor forward one part. At the
last part ($index = of$) it is a no-op. \texttt{T4}: enabled from
$\{ playing, paused \}$. The mode is held, as for \texttt{prev}. This is the
\emph{user-driven} skip; the automatic end-of-part advance is
\texttt{AutoAdvance} below, and unlike it, \texttt{next} at $M$ does not wrap ---
it stalls at the last part.

\begin{schema}{Next}
  \Delta PlayerView
  \where
  mode \in \{ playing, paused \} \\
  catalogued' = catalogued \\
  mode' = mode \\
  album' = album \\
  of' = of \\
  (index < of \implies index' = index + 1) \\
  (index = of \implies index' = index)
\end{schema}

Since $index \leq of$ in every active state, the two cases are exhaustive. When
$index < of$, $index' = index + 1 \leq of$; when $index = of$, $index' = of$;
either way $1 \leq index' \leq of$, so \texttt{T5} holds and the no-op at part
$M$ is the $index = of$ case.

\subsection{AutoAdvance --- the end-of-part advance}
\label{sec:autoadvance}

When the current part ends, the playback loop advances the cursor to another
part of the same album, avoiding an immediate repeat when the album holds more
than one part and replaying when it holds exactly one. This is the phase-2
\texttt{PlayerTrackEnd}, and it is guarded by $advancing = ztrue$ --- equivalently
$mode = playing$ --- so it \emph{cannot} fire while paused. That guard is
\texttt{T3}: a suspended album does not auto-advance.

\begin{schema}{AutoAdvance}
  \Delta PlayerView
  \where
  advancing = ztrue \\
  catalogued' = catalogued \\
  mode' = playing \\
  album' = album \\
  of' = of \\
  1 \leq index' \land index' \leq of \\
  (of \geq 2 \implies index' \neq index) \\
  (of = 1 \implies index' = index)
\end{schema}

The precondition is written $advancing = ztrue$ rather than $mode = playing$ to
make the \texttt{T3} reading literal: the auto-advance fires exactly when the
cursor is progressing, and the state invariant equates that with the playing
mode. The album and its part count are held, so \texttt{T1}, \texttt{T2}, and
\texttt{T7} are preserved; only the cursor moves, and it stays in $1 \upto of$.
Unlike \texttt{Next}, \texttt{AutoAdvance} at the last part does not stall --- it
rotates to another part (or replays the sole part), the existing radio-style
end-of-album behaviour, which is why the two are separate operations and not one.

\section{Design forks surfaced by the formalisation}
\label{sec:forks}

Formalising the transport exposes three points where the brief either departs
from phase~2 or leaves a reading open. They are recorded here for the leader; the
model takes the brief's stated position in each case and flags the tension rather
than papering over it.

\begin{description}
  \item[Fork~A --- play only from idle (\texttt{T4}) versus phase-2's
    start-or-switch. \emph{Resolved: SWITCH.}] The phase-2 \texttt{PlayerPlay}
    (\texttt{docs/vox-music-player.tex}, \S5) had \emph{no} mode precondition: from
    \texttt{idle} it started, while \texttt{playing} it \emph{switched},
    displacing the current album --- the model image of
    \texttt{SwitchSelection}. The phase-3 brief's \texttt{T4} would have
    restricted \texttt{play} to the \texttt{idle} mode. The two cannot both
    hold. The concrete question was: \emph{what happens when a user clicks an
    album-list Play button while another album is playing or paused?} Three
    coherent answers presented themselves: (i) it \emph{switches} --- keep
    phase-2's semantics and drop the $mode = idle$ guard; (ii) it is
    \emph{refused} --- the album-list Play buttons render \texttt{disabled}
    whenever the transport is active, so a switch requires an explicit
    \texttt{stop} first; (iii) it is an implicit \texttt{stop}-then-\texttt{play}.
    \textbf{The leader resolved this to (i), SWITCH.} A normal media player
    switches to a newly clicked album while one is playing, with no stop-first
    required; taking (i) preserves the phase-2 start-or-switch behaviour with no
    regression, and is the simpler design, since the album-list Play buttons
    stay enabled in every mode rather than needing a disabled state. The brief's
    \texttt{T4} ``play only from idle'' clause was too strict and is dropped.
    Accordingly \texttt{PlayAlbum} (Section~\ref{sec:play}) carries \emph{no}
    mode guard: from \texttt{idle} it starts, from \texttt{playing} or
    \texttt{paused} it switches, and in every case the after-state is
    \texttt{target?} playing from part~1.
  \item[Fork~B --- prev/next while paused hold the mode.] The brief says
    prev/next ``move the part cursor within the now-playing album'' and does not
    say they resume. This model holds the mode ($mode' = mode$), so \texttt{prev}
    or \texttt{next} while paused moves the cursor and leaves the player paused.
    The alternative --- prev/next implicitly resume --- would be surprising and is
    not what the brief describes. \textbf{Recommendation:} hold the mode, as
    modelled. Flagged only because ``move the cursor while suspended'' is a state
    the phase-2 player could not reach, and the daemon's \texttt{prev} against a
    \texttt{SIGSTOP}ed subprocess must be well-defined (it must reposition the
    held playback without un-suspending it).
  \item[Fork~C --- next at $M$ stalls; auto-advance at $M$ wraps.] \texttt{Next}
    at the last part is a no-op (it stalls), whereas \texttt{AutoAdvance} at the
    last part rotates to another part or replays the sole part. The brief is
    explicit that these differ (``next at part $M$ is a no-op; auto-advance past
    $M$ is the existing end-of-album behavior''), so this is not a genuine fork,
    only a place the two operations must not be conflated in code: the user's
    \texttt{next} and the loop's end-of-part advance are distinct transitions.
    Recorded so the implementation keeps them apart.
\end{description}

\section{Type-checking and model-checking}
\label{sec:results}

The specification is \texttt{fuzz}-clean: \texttt{fuzz -t} exits~0. It was then
model-checked with ProB (\texttt{probcli}, $maxPart = 3$, the given set
\texttt{ALBUM} at cardinality~2, integers bounded to $0 \upto 3$).

Over the faithful state space --- the boot catalog ranging over every subset of
\texttt{ALBUM}, empty included --- ProB visits every reachable state (53 states,
461 transitions) and finds \emph{no invariant violation}: the \texttt{PlayerView}
predicate, and with it \texttt{T1}, \texttt{T2}, \texttt{T3} (through
\texttt{advancing}), \texttt{T5} (the cursor bounds), \texttt{T6}, and
\texttt{T7}, holds on every reachable state. All seven transitions are covered,
so each is reachable and exercised. \texttt{T4} is enforced by construction: the
five mode-guarded transitions --- \texttt{pause}, \texttt{resume}, \texttt{stop},
\texttt{prev}, \texttt{next} --- each name their enabling mode on the first
predicate line, and no reachable state fires one from a mode the guard forbids;
\texttt{PlayAlbum}, by the Fork~A resolution (Section~\ref{sec:forks}), carries
\emph{no} mode guard and is deliberately enabled from all three modes. ProB
exercises \texttt{PlayAlbum} as a switch from both non-idle modes with no
invariant violation: from a \texttt{playing} state (e.g. \texttt{album} $=
\{$ALBUM2$\}$, replaying or resizing the same source) and from a \texttt{paused}
state (\texttt{glyph} $= playGlyph$), and across albums --- a \texttt{playing}
\texttt{album} $= \{$ALBUM1$\}$ under a two-album catalog steps to
\texttt{album} $= \{$ALBUM2$\}$, \texttt{index} $= 1$, \texttt{mode} $=
playing$, the newly chosen album displacing the active one and restarting at
part~1. The switch adds edges, not states: the after-state of a switch is one a
\texttt{stop}-then-\texttt{play} could already reach, so the state count is
unchanged and only the transition count grows. In particular \texttt{AutoAdvance}
is guarded by $advancing = ztrue$, which the invariant equates with $mode =
playing$, so no reachable \texttt{paused} state enables it --- this is
\texttt{T3}, confirmed by exhaustion.

ProB reports exactly one deadlock, and it is not a defect: the empty-catalog
idle state (\texttt{catalogued} $= \emptyset$, \texttt{mode} $=$ \texttt{idle})
has no enabled transition, because \texttt{PlayAlbum} requires a catalogued
album and this model does not re-state the catalog-mutation operations
(\texttt{music new} / \texttt{CatalogAdd}), which live in phase~2
(\texttt{docs/vox-music-player.tex}) and are unchanged. It is the deliberate
scope boundary of a transport model, not a wedged transport. Restricting the
boot catalog to be non-empty --- so \texttt{PlayAlbum} is always available from
\texttt{idle} --- the machine is \emph{deadlock-free}: ProB visits all~52 states,
fires all~460 transitions, covers all seven operations, and finds no
counterexample of any kind. The transport transitions themselves never wedge;
the only dead end is ``no album has ever been catalogued'', which the phase-2
catalog operations resolve.

\end{document}
